% Generated by scripts/assemble.py; corrections belong in source/corrections.json.
\examhead{42}{上海交通大学 2024 年数学分析试卷}{45}
% p046-b06; source PDF page 46
\par\Needspace{5\baselineskip}
\blockanchor{p046-b06}
\majorquestion{一}判断题. 正确的打“√”，错误的打“×”. 每题6分，共30分.\par
% p046-b07; source PDF page 46
\blockanchor{p046-b07}
\question{1}具有原函数的函数一定 Riemann 可积.\par
% p046-b08; source PDF page 46
\blockanchor{p046-b08}
\question{2}存在 $(- \infty, + \infty)$ 上没一点都不连续的函数.\par
% p046-b09; source PDF page 46
\blockanchor{p046-b09}
\question{3}若 $\int_{a}^{+\infty} f(x) \, \mathrm{d}x$ 收敛, $f(x) > 0$ 且 $f(x)$ 在 $[0, +\infty)$ 上连续, 则 $\lim_{x \to +\infty} f(x) = 0$.\par
% p046-b10; source PDF page 46
\blockanchor{p046-b10}
\question{4}如果函数 $f$ 在 $[a, b]$ 上可导, 而导函数 $f'$ 在 $[a, b]$ 上不连续, 那么 $f'$ 不一定能取到 $f'(a)$ 与 $f'(b)$ 之间的一切值.\par
% p046-b11; source PDF page 46
\blockanchor{p046-b11}
\question{5}级数 $\sum_{n=1}^{\infty} x^{\alpha} e^{-nx} (\alpha > 1)$ 在 $[0, +\infty)$ 上一致收敛。(可能有误)\par
% p046-b12; source PDF page 46
\par\Needspace{5\baselineskip}
\blockanchor{p046-b12}
\majorquestion{二}(25 分) 计算题.\par
% p046-b13; source PDF page 46
\blockanchor{p046-b13}
\question{1}(12 分) 设函数 $f$ 在 $[1, +\infty)$ 上可导, 如果 $\lim_{x \to +\infty} [2f(x) + xf'(x)] = 3$, 求 $\lim_{x \to +\infty} f(x)$.\par
% p046-b14; source PDF page 46
\blockanchor{p046-b14}
\question{2}(13 分) 计算 $\int_{0}^{+\infty} e^{-x^2} \, \mathrm{d}x$.\par
% p046-b15; source PDF page 46
\par\Needspace{5\baselineskip}
\blockanchor{p046-b15}
\majorquestion{三}(25 分) 设 $f(x)$ 为 $(- \infty, +\infty)$ 上的连续函数, 周期函数, $f(x)$ 不为常值函数. 试证明 $f(x)$ 存在最小正周期.\par
% p046-b16; source PDF page 46
\par\Needspace{5\baselineskip}
\par\Needspace{6\baselineskip}
\blockanchor{p046-b16}
\majorquestion{四}(20 分) 求重积分\par
\nopagebreak[4]
% p046-b17; source PDF page 46
\blockanchor{p046-b17}
\displayformula{I = \iint_ {D} \sqrt {\frac {1 - x ^ {2} - y ^ {2}}{1 + x ^ {2} + y ^ {2}}} \mathrm{d} x \mathrm{d} y.}
% p046-b18; source PDF page 46
\blockanchor{p046-b18}
其中 $D$ 是圆周 $x^{2} + y^{2} = 1$ 及坐标轴所围成的第一象限区域.\par
% p046-b19; source PDF page 46
\par\Needspace{5\baselineskip}
\par\Needspace{6\baselineskip}
\blockanchor{p046-b19}
\majorquestion{五}(25 分) 设 $u = u(x_{1}, x_{2}, \cdots, x_{n})$ 是一个二阶连续可导的 $n$ 元函数, $u > 0$, 令\par
\nopagebreak[4]
% p046-b20; source PDF page 46
\blockanchor{p046-b20}
\displayformula{\nabla = \left(\frac {\partial}{\partial x _ {1}}, \frac {\partial}{\partial x _ {2}}, \dots , \frac {\partial}{\partial x _ {n}}\right).}
% p046-b21; source PDF page 46
\par\Needspace{6\baselineskip}
\blockanchor{p046-b21}
$D^2 u = \left(\frac{\partial^2u}{\partial x_i\partial x_j}\right)$ 为矩阵.同时\par
\nopagebreak[4]
% p046-b22; source PDF page 46
\blockanchor{p046-b22}
\displayformula{| \nabla u | = \left(\sum_ {i = 1} ^ {n} u _ {x _ {i}} ^ {2}\right) ^ {\frac {1}{2}}, | D ^ {2} u | = \left(\sum_ {i, j = 1} ^ {n} u _ {x _ {i} x _ {j}} ^ {2}\right) ^ {\frac {1}{2}}.}
% p046-b23; source PDF page 46
\par\Needspace{6\baselineskip}
\blockanchor{p046-b23}
试证明：对任意的 $\theta \in (0,1)$ 成立恒等式\par
\nopagebreak[4]
% p046-b24; source PDF page 46
\blockanchor{p046-b24}
\displayformula{| D ^ {2} u | ^ {2} = \frac {1}{\theta^ {2}} u ^ {2 - 2 \theta} | D ^ {2} u ^ {\theta} | ^ {2} - (1 - \theta) ^ {2} u ^ {- 2} | \nabla u | ^ {4} + 2 (1 - \theta) u ^ {- 1} \nabla u \cdot (D ^ {2} u \cdot \Delta u).}
% p047-b00; source PDF page 47
\par\Needspace{5\baselineskip}
\par\Needspace{6\baselineskip}
\blockanchor{p047-b00}
\majorquestion{六}(25 分) 设 $n \geq 2$ 为整数, 函数 $f \in C^2([0, R])$ 满足 $f_r(0) = 0$. 试证明\par
\nopagebreak[4]
% p047-b01; source PDF page 47
\blockanchor{p047-b01}
\displayformula{\sup _ {r \in (0, R)} \left| \frac {f _ {r} (r)}{r} \right| \leq \frac {1}{n} \sup _ {r \in (0, R)} \left| f _ {r r} + \frac {n - 1}{r} f _ {r} \right|.}
